\documentclass{article}
\begin{document}
Okay, I need to analyze the algorithm that uses online gradients for parameter updates with the goal of achieving sublinear regret. The user provided examples of SGD and Adam, so I should model the answer after those. Let's start by understanding the problem.

First, the algorithm is for online learning where data comes sequentially, and the goal is to minimize regret. Regret is the difference between the cumulative loss of the algorithm and the best fixed point in hindsight. Sublinear regret means that the regret grows slower than the number of steps, like O(√T).

The user wants a formal algorithm description, pseudocode, a dry run example, convergence assumptions, theorems, properties, and a detailed proof with annotated steps.

Starting with the algorithm details. Since it's online with online gradients, it's likely similar to Online Gradient Descent (OGD). OGD updates parameters using the gradient of the current loss function. The formal description would involve an update rule with learning rate η, and gradients computed on incoming data points.

For the mathematical formulation, the update rule would be w_{t+1} = w_t - η_t g_t, where g_t is the gradient of the current function f_t(w_t). The learning rate schedule could be decreasing, maybe η_t = 1/√t to achieve sublinear regret.

Pseudocode: Initialize w_0, loop over time steps, compute gradient, update parameters, record loss.

Dry run example: f(w) = (w - 3)^2, w_0 = 0, η = 0.1. Let's compute gradients for each step. At t=0, w_0=0, gradient is -2*(0-3)=6, so w_1 = 0 - 0.1*6 = -0.6. Then f(w_1) = (-0.6 -3)^2 = (-3.6)^2=12.96. Continue for t=1 and t=2.

Convergence theory: Assumptions would include convexity and Lipschitz continuity of the functions. The main theorem would state that the regret is O(√T). Theoretical properties might compare to other online algorithms, sensitivity to η, etc.

For the proof, I need to follow the reference examples. Start with the regret definition, expand the squared norm, use the update rule, and derive the regret bound. Each step should be annotated.

Wait, the user mentioned "online functions" which are convex and Lipschitz. So each function f_t is convex and L-Lipschitz. The regret is the sum of f_t(w_t) minus the minimum over w of sum f_t(w).

The key steps in the proof would be using the convexity and Lipschitz properties, expanding the update step, and applying the learning rate decay. The regret can be bounded using the difference in squared norms between consecutive iterations.

I need to make sure all steps are detailed, with each non-trivial step annotated. For example, expanding the squared norm, using the update rule, applying the convexity, and summing over iterations.

Let me outline the main proof steps:

1. Define the regret and the goal to bound it.
2. Use the definition of convex functions and the Lipschitz gradient.
3. Expand the squared norm term for each iteration.
4. Apply the update rule to relate g_t to the gradient.
5. Sum over all iterations and telescope the terms.
6. Use the learning rate schedule to bound the sum.
7. Conclude the regret bound.

Also, need to include the variance or unbiasedness if applicable, but since it's online with exact gradients, maybe not necessary. Wait, the problem mentions "online gradients from sequential data", so perhaps exact gradients, not stochastic. So the convergence proof would be deterministic or have different steps compared to SGD.

But in the reference examples, SGD uses stochastic gradients. Here, since it's online with exact gradients (assuming), the proof might be similar to OGD.

I need to structure all this into LaTeX sections as per the user's instructions, without any prose outside the LaTeX. Compilable code only.

Let me start drafting the LaTeX.

First, the algorithm details:

\section*{Online Gradient Descent (OGD)}
\section*{Mathematical Formulation}
The algorithm updates parameters iteratively using the gradient of the current loss function. For convex and L-Lipschitz functions \( f_t(w) \), the update rule is:
\[ w_{t+1} = w_t - \eta_t \nabla f_t(w_t) \]
where \( \eta_t \) is the learning rate at step \( t \).

\section*{Pseudocode}
\begin{algorithm}[H]
  \caption{Online Gradient Descent (OGD)}
  \begin{algorithmic}[1]
    \State Initialize \( w_0 \in \mathbb{R}^d \)
    \For{\( t = 0 \) to \( T-1 \)}
      \State Receive \( f_t \in \{ f_1, f_2, \ldots, f_T \} \)
      \State Compute \( g_t = \nabla f_t(w_t) \)
      \State Update \( w_{t+1} = w_t - \eta_t g_t \)
    \EndFor
  \end{algorithmic}
\end{algorithm}

Dry run example for f(w) = (w-3)^2, w0=0, η=0.1:

Iteration 0:
w0 = 0
g0 = gradient at 0: -2*(0-3) = 6
w1 = 0 - 0.1*6 = -0.6
f(w1) = (-0.6 -3)^2 = 12.96

Iteration 1:
w1 = -0.6
g1 = -2*(-0.6 -3) = -2*(-3.6) = 7.2
w2 = -0.6 - 0.1*7.2 = -0.6 -0.72 = -1.32
f(w2) = (-1.32 -3)^2 = (-4.32)^2 = 18.6624

Iteration 2:
w2 = -1.32
g2 = -2*(-1.32 -3) = -2*(-4.32) = 8.64
w3 = -1.32 -0.1*8.64 = -1.32 -0.864 = -2.184
f(w3) = (-2.184 -3)^2 = (-5.184)^2 ≈ 26.86

Now, moving to the convergence theory:

\section*{Assumptions}
\begin{itemize}
  \item Convexity: \( f_t \) is convex for all \( t \).
  \item Lipschitz continuity: \( \exists L > 0 \) such that \( \|\nabla f_t(w) - \nabla f_t(w')\| \leq L \| w - w' \| \) for all \( w, w' \).
  \item Bounded diameter: \( \exists D > 0 \) such that \( \| w - w' \| \leq D \) for all \( w, w' \in \mathcal{W} \).
  \item Learning rate: \( \eta_t \leq 1/L \) for all \( t \).
\end{itemize}

\section*{Main Convergence Theorem}
The regret \( R(T) = \sum_{t=0}^{T-1} f_t(w_t) - \min_{w \in \mathcal{W}} \sum_{t=0}^{T-1} f_t(w) \) satisfies:
\[ R(T) \leq \frac{D^2 L}{2 \eta} + \frac{\eta L T}{2} \]

\section*{Theoretical Properties}
- The regret is sublinear in \( T \), \( R(T) = O(\sqrt{T}) \) when \( \eta = 1/\sqrt{L T} \).
- Sensitivity to \( L \) and \( D \): Higher Lipschitz constant or diameter increases regret.

\section*{Discussion}
OGD achieves sublinear regret under convex and Lipschitz assumptions. The proof uses the difference of squared norms and learning rate decay. Compared to SGD, it uses exact gradients and has a different regret bound structure.

For the complete proof:

\section*{Preliminaries}
\begin{definition}[Regret]
Regret is \( R(T) = \sum_{t=0}^{T-1} f_t(w_t) - \min_{w \in \mathcal{W}} \sum_{t=0}^{T-1} f_t(w) \).
\end{definition}

\section*{Main Proof}
For each \( t \):
\[
\end{document}